\documentclass[12pt,a4paper]{article}
\usepackage[margin=1in]{geometry}
\usepackage{amsmath,amssymb}
\usepackage{hyperref}
\usepackage{longtable}
\usepackage{xcolor}
\usepackage{listings}
\usepackage{enumitem}
\usepackage[T1]{fontenc}
\usepackage[utf8]{inputenc}

\lstdefinestyle{py}{
  language=Python,
  basicstyle=\ttfamily\small,
  keywordstyle=\color{blue},
  commentstyle=\color{teal},
  stringstyle=\color{red!60!black},
  showstringspaces=false,
  breaklines=true,
  frame=single
}

\title{Deep Learning CT Reconstruction\\Work Summary Report}
\author{Prepared for Project Explanation}
\date{\today}

\begin{document}
\maketitle

\section*{1. What we have done till now}
This project builds a complete computed tomography (CT) reconstruction pipeline and then extends it with a deep-learning based post-processing model. The work can be divided into two major parts:

\begin{enumerate}[leftmargin=*]
\item Classical CT reconstruction from measured projection images.
\item Learning-based enhancement of degraded reconstructions using paired training data.
\end{enumerate}

In the current workspace, the implemented pipeline is:
\begin{enumerate}[leftmargin=*]
\item Load the raw projection TIFF images and scan settings.
\item Parse CT geometry from the \texttt{settings.cto} file.
\item Build an early baseline reconstruction using slice-wise filtered backprojection (FBP).
\item Build the main 3D cone-beam reconstruction using ASTRA FDK.
\item Simulate degraded acquisition conditions such as sparse-view, limited-angle, and noisy scans.
\item Reconstruct both degraded and full-reference volumes.
\item Align and save input-target training pairs.
\item Train a 2D U-Net on axial slices.
\item Export compact training data for Google Colab.
\item Add phantom/STL generation utilities for resolution testing and printable objects.
\end{enumerate}

\paragraph{Important note on chronology.}
The Git history in this repository currently shows only one commit, so this report reconstructs the workflow from the present codebase and the uncommitted files available in the workspace.

\section*{2. Project objective}
The objective is to improve CT image quality when the scan is degraded. In practice, this means:
\begin{itemize}[leftmargin=*]
\item reconstructing a high-quality reference volume from full projection data,
\item creating difficult acquisition scenarios,
\item reconstructing artifact-corrupted volumes from those degraded projections, and
\item training a neural network to map degraded slices to improved slices.
\end{itemize}

Mathematically, if the projection operator is denoted by $\mathcal{P}$ and the ideal object by $f$, then the measured projections are
\[
g = \mathcal{P}f + \epsilon,
\]
where $\epsilon$ represents measurement noise and acquisition imperfections.

The reconstruction operator $\mathcal{R}$ gives
\[
\hat{f} = \mathcal{R}(g).
\]

For degraded acquisition,
\[
g_d = \mathcal{D}(g),
\]
where $\mathcal{D}$ is a degradation operator such as sparse-view sampling, limited-angle selection, or noise addition. Then the degraded reconstruction is
\[
\hat{f}_d = \mathcal{R}(g_d),
\]
and the neural network is trained to learn
\[
\mathcal{N}(\hat{f}_d) \approx \hat{f}_{\text{ref}}.
\]

\section*{3. Step 1: Loading the CT data}
The file \texttt{data\_loader.py} is responsible for reading the dataset from \texttt{sample 1}. It:
\begin{itemize}[leftmargin=*]
\item reads the \texttt{settings.cto} configuration file,
\item loads all TIFF projection images from the \texttt{projections} folder,
\item converts them into a 3D NumPy array of shape
\[
(\text{number of projections}, \text{detector rows}, \text{detector columns}),
\]
\item stores the result in a \texttt{CTScanData} dataclass.
\end{itemize}

The settings parser also converts strings into integers, floating-point values, and booleans where possible.

\begin{lstlisting}[style=py,caption={Core idea of projection loading}]
files = sorted(projections_dir.glob("*.tif"))
stack = np.stack(
    [tifffile.imread(str(path)).astype(dtype, copy=False) for path in files],
    axis=0,
)
\end{lstlisting}

\section*{4. Step 2: Parsing geometry from the scan settings}
The file \texttt{geometry.py} converts the scanner settings into a structured geometry object. This includes:
\begin{itemize}[leftmargin=*]
\item number of projections,
\item total angular range,
\item detector size,
\item detector pixel size,
\item source-to-object distance (SOD),
\item source-to-detector distance (SDD),
\item center of rotation (COR),
\item valid axial reconstruction range \texttt{zmin} to \texttt{zmax}.
\end{itemize}

The projection angles are generated as
\[
\theta_i = \theta_0 + i \cdot \Delta \theta,\qquad i=0,1,\dots,N-1,
\]
where
\[
\Delta \theta = \frac{\text{angle range}}{N}.
\]

In code, the angles are created using:
\begin{lstlisting}[style=py,caption={Angle generation}]
angles_deg = np.linspace(start_angle_deg, stop_angle_deg,
                         projections, endpoint=False, dtype=np.float32)
angles_rad = np.deg2rad(angles_deg).astype(np.float32)
\end{lstlisting}

\section*{5. Step 3: First baseline approach using slice-wise FBP}
The file \texttt{reconstruct\_fbp.py} implements an early baseline. This is not the final cone-beam solution, but it is important because it establishes the basic CT reconstruction idea.

\subsection*{5.1 Sinogram construction}
For a chosen detector row, the code extracts one horizontal line from each projection. This produces a 2D sinogram:
\[
S(s,\theta),
\]
where $s$ is detector position and $\theta$ is projection angle.

\begin{lstlisting}[style=py,caption={Constructing a sinogram from a detector row}]
def construct_sinogram(projections, detector_row):
    return projections[:, detector_row, :].T
\end{lstlisting}

\subsection*{5.2 Ramp filtering}
Classical FBP requires filtering in frequency space. If $\mathcal{F}$ is the Fourier transform, then the filtered sinogram is
\[
\tilde{S}(s,\theta) = \mathcal{F}^{-1}\left( |\omega| \, \mathcal{F}(S(s,\theta)) \right),
\]
where $|\omega|$ is the ramp filter.

This is implemented as:
\begin{lstlisting}[style=py,caption={Ramp filter implementation}]
frequencies = np.fft.fftfreq(num_detectors, d=detector_spacing_mm)
ramp = np.abs(frequencies)
spectrum = np.fft.fft(sinogram, axis=0)
filtered = np.fft.ifft(spectrum * ramp[:, None], axis=0).real
\end{lstlisting}

\subsection*{5.3 Backprojection}
For each image point $(x,y)$ and each angle $\theta$, the detector coordinate is approximated by
\[
s = x\cos\theta + y\sin\theta.
\]

The reconstruction is then accumulated over all angles:
\[
f(x,y) \approx \int_0^\pi \tilde{S}(x\cos\theta + y\sin\theta,\theta)\, d\theta.
\]

The code performs this discretely:
\begin{lstlisting}[style=py,caption={Discrete backprojection loop}]
for angle_index, theta in enumerate(angles_rad):
    detector_coordinate = xx * np.cos(theta) + yy * np.sin(theta)
    sampled = np.interp(detector_coordinate.ravel(),
                        detector_positions_mm,
                        filtered[:, angle_index],
                        left=0.0, right=0.0)
    reconstruction += sampled.reshape(image_size, image_size)
\end{lstlisting}

\subsection*{5.4 Why this baseline matters}
This baseline helped us:
\begin{itemize}[leftmargin=*]
\item verify that the data loading and angle handling were correct,
\item understand the basic sinogram-to-image relationship,
\item produce an interpretable first reconstruction before moving to 3D cone-beam reconstruction.
\end{itemize}

\section*{6. Step 4: Main 3D reconstruction using ASTRA FDK}
The file \texttt{reconstruct\_fdk\_astra.py} is the main reconstruction script. It uses the ASTRA toolbox with the CUDA FDK algorithm for cone-beam CT.

\subsection*{6.1 Detector downsampling}
Because full-resolution detector data can exceed GPU memory, the code spatially downsamples each projection by block averaging:
\[
P_{\text{ds}}(i,j) = \frac{1}{k^2}\sum_{u=0}^{k-1}\sum_{v=0}^{k-1} P(ki+u, kj+v).
\]

\begin{lstlisting}[style=py,caption={Spatial projection downsampling}]
return image.reshape(cropped_h // factor, factor,
                     cropped_w // factor, factor).mean(axis=(1, 3))
\end{lstlisting}

\subsection*{6.2 Intensity to attenuation conversion}
CT reconstruction should use line integrals of attenuation. If $I$ is the measured intensity and $I_0$ is the air intensity, then:
\[
\mu L = -\log\left(\frac{I}{I_0}\right).
\]

The code estimates $I_0$ from a high percentile of each projection:
\begin{lstlisting}[style=py,caption={Beer-Lambert conversion}]
air_level = np.percentile(projections, air_percentile, axis=(1, 2), keepdims=True)
normalized = np.clip(projections / air_level, 1e-6, 1.0)
attenuation = -np.log(normalized)
\end{lstlisting}

\subsection*{6.3 Cone-beam geometry}
The ASTRA projection geometry is created using:
\begin{itemize}[leftmargin=*]
\item detector pixel spacing,
\item detector rows and columns,
\item projection angles,
\item source-to-object distance,
\item source-to-detector distance.
\end{itemize}

The volume voxel size is estimated by geometric magnification:
\[
\text{voxel size} = \text{detector pixel size} \cdot \frac{\text{SOD}}{\text{SDD}}.
\]

\subsection*{6.4 FDK reconstruction}
FDK is a cone-beam extension of filtered backprojection. In simplified form, it:
\begin{enumerate}[leftmargin=*]
\item applies cone-beam weighting,
\item filters projections,
\item backprojects them through 3D geometry.
\end{enumerate}

The code delegates this to ASTRA:
\begin{lstlisting}[style=py,caption={Running FDK with ASTRA}]
cfg = astra.astra_dict("FDK_CUDA")
cfg["ProjectionDataId"] = proj_id
cfg["ReconstructionDataId"] = vol_id
alg_id = astra.algorithm.create(cfg)
astra.algorithm.run(alg_id)
reconstructed = astra.data3d.get(vol_id).astype(np.float32, copy=False)
\end{lstlisting}

\subsection*{6.5 Cropping valid z-range}
After reconstruction, only the physically valid axial range is kept:
\[
z \in [z_{\min}, z_{\max}].
\]
This removes unsupported slices outside the meaningful reconstruction interval.

\section*{7. Step 5: Simulating degraded acquisition conditions}
The file \texttt{simulate\_degradation.py} generates harder CT scenarios so that the deep model learns to correct artifacts.

\subsection*{7.1 Full dataset}
The full dataset is used as the reference case. It keeps all projections and all angles.

\subsection*{7.2 Sparse-view CT}
Sparse-view acquisition keeps every $k$-th projection:
\[
g_{\text{sparse}} = \{g_0, g_k, g_{2k}, \dots\}.
\]

This reduces the number of views and typically produces streak artifacts.

\begin{lstlisting}[style=py,caption={Sparse-view selection}]
indices = np.arange(0, projections.shape[0], step, dtype=np.int32)
return projections[indices], angles_rad[indices], indices
\end{lstlisting}

\subsection*{7.3 Limited-angle CT}
Limited-angle acquisition keeps only a restricted angular interval:
\[
\theta \in [\theta_{\text{start}}, \theta_{\text{stop}}].
\]

This makes the inverse problem incomplete and usually introduces directional artifacts and missing-edge effects.

\subsection*{7.4 Noise simulation}
Two noise models are implemented.

\paragraph{Gaussian noise.}
\[
g_{\text{noisy}} = g + \eta,\qquad \eta \sim \mathcal{N}(0,\sigma^2).
\]

\paragraph{Poisson noise.}
Photon-counting physics is approximated by:
\[
c \sim \text{Poisson}(\lambda),
\]
followed by a scaling back to intensity space. This better reflects shot noise in X-ray imaging.

\begin{lstlisting}[style=py,caption={Poisson noise simulation}]
counts = rng.poisson(scaled * photon_budget).astype(np.float32)
restored = np.clip(counts / photon_budget, 1e-6, None) * peak
\end{lstlisting}

\subsection*{7.5 Default degraded sets}
The code currently prepares these standard cases:
\begin{itemize}[leftmargin=*]
\item sparse-view with step 2,
\item sparse-view with step 4,
\item limited-angle from $0^\circ$ to $180^\circ$,
\item limited-angle from $30^\circ$ to $210^\circ$,
\item Poisson noise with level 0.5,
\item Poisson noise with level 0.25.
\end{itemize}

\section*{8. Step 6: Building supervised training pairs}
The file \texttt{build\_training\_pairs.py} creates paired data for learning.

\subsection*{8.1 Reference target volume}
First, the code reconstructs a reference volume from the full dataset using the FDK pipeline.

\subsection*{8.2 Input volume}
Then it reconstructs each degraded dataset with the same FDK pipeline. This ensures that the learning task compares:
\[
\text{degraded FDK reconstruction} \rightarrow \text{reference FDK reconstruction}.
\]

\subsection*{8.3 Shape alignment}
Because two reconstructions may differ slightly in size after downsampling and cropping, the code performs center-cropping to their common shape:
\[
(Z,Y,X)_{\text{shared}} = \min\big((Z,Y,X)_{\text{input}}, (Z,Y,X)_{\text{target}}\big).
\]

\begin{lstlisting}[style=py,caption={Shape alignment by center crop}]
shared_shape = tuple(min(a, b) for a, b in zip(input_volume.shape, target_volume.shape))
\end{lstlisting}

\subsection*{8.4 Slice extraction and normalization}
The volumes are normalized to the range $[0,1]$:
\[
v_{\text{norm}} = \frac{v - v_{\min}}{v_{\max} - v_{\min}}.
\]

Then axial slices are saved into a compressed \texttt{.npz} file:
\begin{itemize}[leftmargin=*]
\item \texttt{input\_slices}
\item \texttt{target\_slices}
\end{itemize}

This prepares a supervised learning dataset at the slice level.

\section*{9. Step 7: Training the 2D U-Net}
The file \texttt{train\_unet.py} trains a 2D U-Net on axial slice pairs.

\subsection*{9.1 Why 2D slices}
Instead of training a full 3D network immediately, the code uses 2D axial slices because:
\begin{itemize}[leftmargin=*]
\item memory cost is lower,
\item training is easier to start and debug,
\item data preparation is simpler,
\item it is a reasonable first deep-learning baseline.
\end{itemize}

\subsection*{9.2 Dataset format}
Each training item is:
\[
x \in \mathbb{R}^{1 \times H \times W},\qquad y \in \mathbb{R}^{1 \times H \times W},
\]
where $x$ is the degraded axial slice and $y$ is the reference slice.

\subsection*{9.3 Network structure}
The network contains:
\begin{itemize}[leftmargin=*]
\item encoder blocks with two convolutions each,
\item max-pooling for downsampling,
\item a bottleneck,
\item transpose convolutions for upsampling,
\item skip connections between encoder and decoder.
\end{itemize}

This architecture is well suited for image restoration because it preserves both global context and local detail.

\begin{lstlisting}[style=py,caption={U-Net block idea}]
self.net = nn.Sequential(
    nn.Conv2d(in_channels, out_channels, kernel_size=3, padding=1),
    nn.ReLU(inplace=True),
    nn.Conv2d(out_channels, out_channels, kernel_size=3, padding=1),
    nn.ReLU(inplace=True),
)
\end{lstlisting}

\subsection*{9.4 Loss function}
The model uses L1 loss:
\[
\mathcal{L}_{L1} = \frac{1}{N}\sum_{i=1}^N |\hat{y}_i - y_i|.
\]

This is a common choice for reconstruction and denoising because it often preserves edges better than mean squared error.

\subsection*{9.5 Optimization}
The optimizer is Adam with configurable learning rate. The training loop performs:
\begin{enumerate}[leftmargin=*]
\item forward pass,
\item loss computation,
\item backpropagation,
\item parameter update.
\end{enumerate}

\begin{lstlisting}[style=py,caption={Training step}]
optimizer.zero_grad(set_to_none=True)
predictions = model(inputs)
loss = criterion(predictions, targets)
loss.backward()
optimizer.step()
\end{lstlisting}

\subsection*{9.6 Validation metric}
The code also computes PSNR:
\[
\text{PSNR} = 10 \log_{10}\left(\frac{1}{\text{MSE}}\right),
\]
assuming the normalized intensity range is $[0,1]$.

Higher PSNR generally indicates better agreement with the target slices.

\section*{10. Step 8: Exporting data for Colab}
The file \texttt{export\_for\_colab.py} copies only the compact files needed for training:
\begin{itemize}[leftmargin=*]
\item \texttt{axial\_slices.npz}
\item \texttt{pair\_metadata.json}
\end{itemize}

Then it creates a manifest and an optional ZIP archive. This makes it easy to train the model in Google Colab without transferring all reconstruction outputs.

\section*{11. Step 9: Phantom and STL generation work}
The current workspace also contains newer scripts that are not committed to Git yet. These scripts extend the project beyond reconstruction into test-object design and fabrication.

\subsection*{11.1 USAF-like phantom generator}
The file \texttt{usaf\_phantom\_generator.py} generates a 3D phantom based on USAF resolution elements. It computes line-pair frequency as:
\[
\text{lp/mm} = 2^{g + (e-1)/6},
\]
where $g$ is the group and $e$ is the element number.

Then:
\[
\text{line width} = \frac{1}{2\cdot \text{lp/mm}},\qquad
\text{pitch} = \frac{1}{\text{lp/mm}}.
\]

These patterns are rasterized into 2D bar structures and extruded into a 3D volume.

\subsection*{11.2 Printable USAF STL}
The file \texttt{export\_usaf\_phantom\_stl.py} builds a classical MIL-style printable resolution plate using repeated rectangular boxes. Each bar is represented as a cuboid in 3D and then exported as ASCII STL.

\subsection*{11.3 Resolution cuboid with internal features}
The file \texttt{export\_resolution\_cuboid\_stl.py} creates a cuboid containing enclosed internal bar-space patterns. This is useful for testing whether CT can resolve hidden internal features of different sizes.

The internal structure is represented as a voxel occupancy grid, and exposed surfaces are converted to triangles for STL output.

\subsection*{11.4 STL scaling utility}
The file \texttt{scale\_ascii\_stl.py} rescales ASCII STL geometry by multiplying each vertex coordinate by a chosen scale factor:
\[
(x,y,z) \mapsto (\alpha x, \alpha y, \alpha z).
\]

\section*{12. Summary of all approaches we implemented}
Till now, the project includes the following technical approaches:

\begin{longtable}{|p{3.1cm}|p{4.0cm}|p{7.2cm}|}
\hline
\textbf{Approach} & \textbf{File(s)} & \textbf{Purpose} \\
\hline
Slice-wise FBP baseline & \texttt{reconstruct\_fbp.py} & Early classical reconstruction baseline for understanding sinograms, filtering, and backprojection. \\
\hline
3D cone-beam FDK & \texttt{reconstruct\_fdk\_astra.py} & Main reconstruction pipeline using ASTRA CUDA. \\
\hline
Degradation simulation & \texttt{simulate\_degradation.py} & Create sparse-view, limited-angle, and noisy projection datasets. \\
\hline
Supervised pair creation & \texttt{build\_training\_pairs.py} & Produce aligned degraded/reference volumes and normalized axial slices. \\
\hline
Deep learning restoration & \texttt{train\_unet.py} & Train a 2D U-Net to map degraded slices to clean slices. \\
\hline
Portable export workflow & \texttt{export\_for\_colab.py} & Export compact training data for cloud training. \\
\hline
Resolution phantom generation & \texttt{usaf\_phantom\_generator.py} & Create a 3D USAF-like phantom volume. \\
\hline
Printable phantom STL & \texttt{export\_usaf\_phantom\_stl.py} & Build a printable USAF-style test plate. \\
\hline
Internal resolution cuboid & \texttt{export\_resolution\_cuboid\_stl.py} & Build a voxel-based STL with enclosed test patterns. \\
\hline
STL scaling utility & \texttt{scale\_ascii\_stl.py} & Uniformly scale printable geometry. \\
\hline
\end{longtable}

\section*{13. Overall workflow in simple words}
If I explain the project simply to a professor, I would say:

\begin{quote}
We first loaded raw CT projection images and scanner geometry. Then we tried a basic reconstruction method using filtered backprojection on selected detector rows to understand the reconstruction process. After that, we implemented a full 3D cone-beam reconstruction using the ASTRA toolbox and the FDK algorithm. Next, we artificially degraded the scan data using sparse angles, limited angle ranges, and noise to simulate difficult real-world CT conditions. We reconstructed both the clean and degraded data, paired them, and trained a 2D U-Net model so that it learns to improve degraded reconstructions slice by slice. In addition, we started generating phantom and STL models for resolution analysis and possible physical testing.
\end{quote}

\section*{14. Code snippets representing the main pipeline}

\subsection*{14.1 Build full reference reconstruction}
\begin{lstlisting}[style=py]
dataset = make_full_projection_dataset(sample_dir)
volume, info = reconstruct_volume_from_projection_dataset(
    dataset=dataset,
    sample_dir=sample_dir,
    downsample_factor=2,
)
\end{lstlisting}

\subsection*{14.2 Build degraded datasets}
\begin{lstlisting}[style=py]
datasets = [
    create_sparse_view_dataset(sample_dir=sample_dir, step=2),
    create_sparse_view_dataset(sample_dir=sample_dir, step=4),
    create_limited_angle_dataset(sample_dir=sample_dir, start_deg=0.0, stop_deg=180.0),
    create_noisy_dataset(sample_dir=sample_dir, mode="poisson", level=0.25, seed=0),
]
\end{lstlisting}

\subsection*{14.3 Build training pairs}
\begin{lstlisting}[style=py]
outputs = build_default_training_pairs(
    sample_dir="sample 1",
    downsample_factor=2,
    output_dir="outputs/training_pairs",
)
\end{lstlisting}

\subsection*{14.4 Train the model}
\begin{lstlisting}[style=py]
python train_unet.py \
  --pairs-root outputs/training_pairs \
  --output-dir outputs/unet_training \
  --epochs 20 \
  --batch-size 8 \
  --learning-rate 1e-3
\end{lstlisting}

\section*{15. What can be said as conclusion}
Till now, the work shows a complete research workflow:
\begin{itemize}[leftmargin=*]
\item raw CT data understanding,
\item classical reconstruction,
\item advanced 3D cone-beam reconstruction,
\item artifact simulation,
\item supervised data generation,
\item deep-learning enhancement,
\item phantom generation for further testing.
\end{itemize}

So this is not just one script. It is a full pipeline from acquisition data to reconstruction to machine-learning based quality improvement.

\section*{16. Possible next extensions}
Logical next steps for the project are:
\begin{itemize}[leftmargin=*]
\item quantitative evaluation using PSNR, SSIM, and error maps on saved outputs,
\item inference script for applying the trained U-Net to whole volumes,
\item 3D deep-learning models instead of 2D slices,
\item comparison with iterative reconstruction methods,
\item evaluation of the generated phantoms with the same CT pipeline.
\end{itemize}

\end{document}
